\chapter{Fazit}\label{ch:fazit}

Ziel dieser Arbeit war es, einen unternehmensinternen Chatbot mit integrierten Tool Calling auf Schwachstellen gegenüber indirekten Prompt-Injection-Angriffen zu 
untersuchen. Anhand der festgesellten Schwachstellen sollte ein Konzept für eine Schutzmaßnahme entwickelt werden, welche ohne Änderungen an dem System integriert 
werden kann. Diese Maßnahme wurde implementiert und die Werte anhand von Metriken gemessen. Die gesammlten Ergebnisse wurden in Kapite~\ref{ch:evaluation} 
ausgewertet. In diesem Kapitel sollen ur mithilfe dieser Evaluationen die Forschungsfrage mit den Teilforschungsfragen beantwortet werden. Daraufhin wird eine 
Hadnlungsempfehlung für den Chatbot4U definiert und ein Ausblick gegeben. 

\section{Beantwortung der Forschungsfragen}

\paragraph{Teilforschungsfrage 1: Welche Schwachstellen weist der Chatbot in Bezug auf indirekte Prompt-Injection-Angriffe auf?}

Aus den in Kapite~\ref{ch:methodik} durchgeführten Angriffen wurde deutlich, wie anfällig der Chatbot auf indirekte Prompt-Injection-Angriffe war. In der lokalen
Testumgebung waren von sieben durchgeführten Angriffen sechs in mindestens nem Versuch erfolgreich. Der Chatbot verhinderte nicht die Ausführung versteckten Inhalts,
das Übermitteln von Daten an eine lokale Wbsebsiete und folgte Aufforderungen den Nutzer über die Aktionen nicht zu informationene. Die einzelnen Versuche der
Angriffe wurden jeweils nur einmal durchgeführt, weshalb sich keine Aussag darüber treffen lässt, mit welcher Wahrscheinlihckeit diese Angriffe vom \ac{llm} 
zugelassen werden. \\ 
Es konnten jedoch anhand dieser erfolgreich ausgeführten Angriffe fünf Schwachstellen definiert werden:

\begin{enumerate}
    \item Fehlende Trennung von Anweisungen und Daten
    \item Übermäßige Funktionalität
    \item Fehlende Kontrolle des Datenabflusses
    \item Keine Vorverarbeitung externer Inhalte
    \item Eingeschränkte Nachvollziehbarkeit
\end{enumerate}

Der Chatbot erfüllte durch seine Fähigket Tool Calls durchzuführen alle drei Kriterien des Lethal Trifecta. Er erlaubte den Zugriff auf private und möglicherweise 
sensible Daten, die Verarbeitung von nicht vertrauenswürdigen Inhalten und die Möglichkeit eine externe Kommunikation durchzuführen. 

\paragraph{Teilforschungsfrage 2: Wie zuverlässig ist ein vorgeschaltetes LLM bei der Abwehr von Prompt-Injection-Angriffen?}

Bei der Verhinderung von Prompt-Injection-Angriffen ist ein vorgeschaltetes \ac{llm} je nach Design weniger hilfreich. Bei einer einfachen binären Abfrage mit dem 
\ac{llm} Mistral Small 4 weden bei einem eigenen Datensatz 73,1\,\% der Angriffe abgewehrt. Bei einem Datensatz aus der Kombination von BIPIA und NotInject liegt 
dieser Wert nur bei 34,0\,\% und bei einem verschleierten BIPIA Datensatz bei 43,0\,\%. Bei der Verwendung von dem \ac{llm} Mistral Large 4 wurden dahingegen 
bei dem eigenen Datensatz 41,6\,\% und bei dem Datensatz BIPIA und NotInject 24,5\,\% der Angriffe erkannt. Gleichzeitig wurde die \ac{fpr} von 1,2\,\% erreicht. \\
Diese Ergebnisse sind Abhängig von dem gestellten Prompt, welcher dem \ac{llm}
mit der Nutzeranfrage mitgeliefert wird, sowie der schädlichen Anweisungen, die das \ac{llm} finden muss und der Art des Inhalts. Dazu kommen die Auswirkungen, welche
durch Triggerwörter entstehen, die durch das \ac{llm} voreilig als Angriff gewertet werden. \\
Somit wird bei der Beantwortung der Teilforschungsfrage klar, dass ein vorgeschaltete \ac{llm} bei der Erkennung von indirekten Prompt-Injection Angriffen nur bedingt
erfolgreich ist und je komplexer ein ANgriff ist, desto schlechter wird die Leistung. Zusätzlich wurde die Frage für den untersuchten Ansatz und nur \acp{llm} von 
Mistral beanwortet und somit sehr nur für einen sehr spezifischen Anwendungsfall beantwortet.

\paragraph{Teilforschungsfrage 3: Verbessert sich die Erkennung von Prompt-Injection-Angriffen, wenn die Schutzmaßnahme um Segmentierung und Rekonstruktion erweitert wird?}

Durch die Erweiterung der Schutzmaßnahme konnten starke Veränderungen festgestellt werden. Der Recall bei der Erkennung von ANgriffen bei dem eigenen Datensatz stieg
von 73,1\,\% auf 82,8\,\%. Bei dem Datensatz BIPIA und NotInject stieg es von 34,0\,\% auf 92,0\,\% und bei dem verschleierten BIPIA stieg es von
43,0\,\% auf 91,0\,\%. Auswirkungen auf die \ac{fpr} waren nicht so stark und blieben bei den öffentlichen Datensätzen identisch und sank bei dem eoigenen Datensatz
von 10,9\,\% auf 6,8\,\%. \\
Die unterschiedlichen Stufen der weiterentwicklung des NAsatzes haben verschiedene Auswikrungen bei dieser Verbesserung verursacht. Die Rekonstruktion, welche in 
der ersten Stufe des Ansatzes durchgeführt wird, ist vor allem hilfreich bei verschleierten Angriffen aus dem eigenen Datensatz und dem verschleierten BIPIA
Datensatz. Durch die Entschlüsselung von 
Base64 kodierten Anfragen und einer Rechtschreibkorrektur und weiteren rekonstruktionen können Angriffe beser klassifiziert werden. Mit der Verwendung von Leitfragen, 
um die Absicht einer Anfrage zu identifizieren konnten vor allem bei dem eigenen Datensatz die Anzhal der richtig klassifizierten Angriffe stark erhöht werden. Diese,
in Stufe 2, udurchgeführte Maßnahme trägt wesenlich dazu bei Angriffe zu erkennen. Durch die Segmentierung, welche in der dritten Stufe der Maßnahme durchgeführt wird,
aber nur bei externen Inhalten wird es dem \ac{llm} ermöglicht eine besser Übersicht bei langen Daten zu behalten. Dadurch das alle Zeilen nummeriert werden und einzelne 
behnaldet werden, ist es nicht mehr ausschlaggebend, an welcher Position sich der ANgriff befindet und wie lang der externe Inhalt ist. Dies ist vor allem für den
BIPIA Datensatz von Bedeutung, bei dem 86,5\,\% der ANgriffe dadurch erlannt werden. \\
Für die teilfprschungfrage lässt sich somit die Aussage treffen, dass ie Verbeserung hilfreich, bei der Erkennung von Angriffen ist. Damit zusammenhängend ist jeodch
eine starke erhöhung der notwendigen Anfragen, der Kosten und der Latenz. Es können bis zu 6 Anfragen notwendig sein, Die Kosten können bis um das 12-Fache steigen 
und die Latenz steigt auf 6\,s bei komplexen Anfragen.

\paragraph{Teilforschungsfrage 4: Inwieweit lassen sich die Ergebnisse auf öffentliche Datensätze übertragen?}

Bei den gesammelten Ergebnissen wird deutlich, dass sich die Leistungsfähigkeit der Erknnung von ANgriffen bei den öffentlichen Datensätzen in Bezug auf den 
Originalansatz am stärksten angestiegen sind. Die Werte stiegen beim Datensatz BIPIA und NotInject von 34,0\,\% auf 92,0\,\% und von 43,0\,\% auf 91,0\,\% bei 
dem verschleierten BIPIA Datensatz. Es wurden jedoch die Schwächen festgestellt, dass Fehlalarme stark von der Art der Daten abhängt. besonders viele 
Fehlalamre entstanden dabei durch Codeausschnite oder E-Malis. Zusätzlich konnte festgestellt werden, dass viele der Fehlalarme welch beim Originalansatz enstanden 
durch die weiterentwicklung des Ansatzes nicht behoben werden konnten. \\
Die Frage kann beantwortet werden, dass er zwar übertragbar ist, jedoch nur beschränkt, da zum ienen die verwendeten öffentlichen Datensätze auf Englisch waren und 
dadurch für das Tsten von externen Inhalten nur Egnlischsprache Datensatzeintrage verwendet wurden, obwohl der Ansatz auf deutsch formuliert wurde. Zusätzlich wurde
der Ansatz nur für tabellen, Codeausschnitte und E-Mails getestet, welches die übertragbarkeit ebenfalls einschränkt. 

\paragraph{Hauptforschungsfrage: Inwieweit kann ein unternehmensinterner Chatbot mit implementiertem Tool Calling gegen indirekte Prompt-Injection-Angriffe abgesichert werden?}

Mit den gesammelten Ergebnissen kann man die Aussage treffen, dass ein vorgeschaltetes \ac{llm} einen unternehmensinternen Chatbot gegen indirekte Prompt-Injection-
Angriffe zum teil absichenrm kann. Es werden mit mit dem Ansatz statt 34,0\,\% bis 73,1\,\% der ANgriffe 83\,\% bis 92\,\% richtig als Angriff klassifiziert. Bei 
den aus Kapitel~\ref{ch:methodik} durchgeführten Angriffsvarianten konnten mit dem ANsatz 35 von 37 richtig erkannt werden. Jedoch wurde der Ansatz nicht in den Chatbot
integriert sondern unabhängig getstet. Somit ist ist die ASussagekraft dahingegen eingeschränkt. \\
Der Ansatz versagt an aussagekraft, da zum einen in einem realiltäts nahen betrieb des Chatbots im Kundenservice, die erreichte \ac{fpr} den Betrieb behindern würde. 
Wenn von den Kundenanfragen 7 bis 14\,\% fälschlicherweis als ANgriff klassifiziert wird und abgeehnt wird, würde dies den Service deutlich schädigen. Zusätzlich sind
keine Multi-turn ANgriffe erkennbar, da die Prüfung einer Anfrage bei jeder unabhängig ist und keine aufeinander aufbauenden vorgesheen sind. AUch das die erkennung,
ob es sich um externen Inhalt handelt wird vorab im Datensatz defineirt und kann mit der SChutzmaßnahme nicht überprüft werden. Dies ist ebenfalss für die 
praktische Anwendung unpraktisch. \\
Die Schwachstellen des Chatbots, wie die übermäßige Funktionalität, wurden daher nicht direkt behoben, sondern mithilfe der ausfürhlichen Kontrolle der 
Nutzeeingabe, die AUswirkuge eingeschränkt, welche durch die Verwendung von Tools entstehen könnte.
Ein vorgeschaltetes \ac{llm} kann die Erkennung deutlich verbessern, bietet allein aber keine ausreichende Absicherung. Diese ist nur in Kombination mit 
systemseitigen Maßnahmen erreichbar.

\section{Handlungsempfehlung}

Bei dem in dieser Arbeit untersuchten System ChatBot4U wurden in dem Abschnitt~\ref{sec:schwachstellen_umgebung} anhand von durchgeführten ANgriffen die 
Schwachstellen identifiziert. Es wurde anhand dieser Schwachstellen versucht eine Schutzmaßnahme zu entwickeln, welche den Chatbot absichert und die Schwachstellen
eindämmt. Mit der Evaluation in Kaptel~\ref{ch:evaluation} wird jedoch deutlich, dass dies nicht vollständig möglich war. Zusätzlich wurde die Maßnahme nicht in 
den Chatbot integriert, sondern unabhängig bewertet, wodurch zusätzlich über manche der Schwachstellen keine Aussage getroffen werden konnte, ob diese durch die
Maßnahme eingedämmt werden konnten oder nicht. Mit den Ergebenisse können nun mehrere Handlungsempfehlungen vorgestellt werden, die die Eindämmung der Schwachstellen
untersützt. Es wurden insgesamt vier entwickelt. \\

\paragraph{Einschränkung übermäßiger Funktionalität}

Einem \ac{llm} ist es möglich auf jegliche Zielwebseiten zuzugreifen und Aktionen auszuüfhren. Es sind alle drei Vorraussetzung des Lethal Trifecta füllt. Eine 
Möglichkeit dies einzuschränken ist, das die Funktion \texttt{auditUrl} angepasst wird. Bisher wird nur üborpft, ob die Zieladress lokal oder öffentlich 
verfügbar ist und ob es sich um einen POST oder GET Anfrage handelt. Jedoch gibt es trotz der Überprüfung keine Einschänkungen bei den ZUgriffen. Aus diesme 
Grund ist es notwendig die Funktionalität, dass Aktionen auf allen Wbsiten und der Zugriff möglich sind eingeschränkt werden. Die einfachste Möglichkeit ist dabei
eine Liste von Ziladressen zu integrieren, mit denen das \ac{llm} in seinem Zugriff eingeschränkt wird.

\paragraph{Menschliche Kontrolle des Datenabflusses}

Ausschlaggebenede Schwachstellen, welche fetgesetellt werden konnten sind, dass es einem ANgreifer möglich war auf eine Zielwebseite zuzugreifen und Daten an diese
zu senden, ohne ein Nutzer die Aktion nachvollziehen konnte. Es sollte daher eine menschliche Kontrolle integriert werden mit der der Nutzer vor der Verwendung
eines Tools eine Beschreibung erhält, welche Aktion ausgeführt werden soll, welches Tool gewählt wurde und welche Daten übergeben werden sollen. Damit kann 
verhindert werden, dass ein \ac{llm} den Nutzer täuscht und ein anderen Tool als notwendig ausführt.

\paragraph{Erkennung externer Inhalte}

Bei dem in dem untersuchten System integrierten Tool Calling wir bei der ausführung eines Aufrufes einer Webseite der gesamte Inhalt als \ac{html} Quelltext übergeben.
Zusätzlich ist in der Schutzmaßnahme keine Prüfung auf externn Inhalt integriert, sondern es wird vorher im Datensatz fest definiert. Extrene Inhalte sollten 
bevor sie Überprüft werden aussortiert werden und aktiv als externer Inhalt gekennzeichnet werden. Dadurch wäre es unter anderem möglich, dass die \acp{llm} nicht 
davon beeinflusst werden, dass sich beispielweise versteckerter Inhalt oder unsichtbare Zeichen in der Anfrgae befinden, welcher Anweisungen beinhalten könnte.

\paragraph{Implementierung der Schutzmaßnahme}

Die in dieser Arbeit entwickelte Schutzmaßnahme kann in Chatbots integriert, um den Schutz zu verstärken. Jedoch sollte beachtete, dass er nicht die einzige 
Schutzmaßnahme ist und auch nicht für den gesamten Verkehr eingesetzt wird. Bei den Ergebnisse wurde deutlich, dass der Ansatz mit einer \ac{fpr} vin bis zu 
13,5\,\% und einer Latenz von bis zu 6\,s der Einsatz im Kundenservice zu Problemen führen könnte. Die Benutzerfreundlichkeit für den Nutzer für reduzuert werden
und da eine Anfrage, wenn ein Angriff erkannt wird, abgelehnt wird, ohne das ein Nutzer eine Begündung erhält würde dies zu Unzufriedenheit führen. Es empfiehlt
sich also die Schutzmaßnahme für die Überprüfung externer Inhalte zu integrieren und anstatt dass eine Anfrage ohne Begründung abgelehnt wird, dem Haupt-\ac{llm}
die Anfrage zusammen mit dem Verdacht zu überreichen und dieses die endgültige Entscheidung treffen zu lassen. Zusätzlich sollte bei der Integration der
SChutzmaßnahme die Anfragen parallelisiert werden, um die Latenz zu reduzieren.

\section{Ausblick}

Diese Arbeit und die gesammelten Ergebniss liefert einen Ausfürhlichen Eindruck, welche Ansätze notwendig sind bei einer weiterführenden Untersuchen.\\
Der wichtigste Punkt ist dabei, dass der Ansatz in die produktiv Umgebung implementiert wird. Damit kann ein Eindruck der Leistung in einer praxisnahen Anwendung 
enstehen und die Einsatzmöglichkeiten besser beurteilt werden. \\
Die Schutzmaßnahme sollte außerdem mit verschiedenn \acp{llm} getestet werden. Es bleibt ungeklärt, ob andere Modelle mit dem Originalansatz in besserer Ergebnis
liefern und somit die Schutzma0nahme irrelevant ist. Díes ist auch in Bezug auf die Kosten von bedeutung, da eine Püfung einer Anfrage kostentechnisch fast gleich 
auf ist mit einer Anfrage an das \ac{llm} Mistral Large 4. \\
Die Maßnahme arbeitet bis jetzt ausschließlich mit einer binären Entscheidung. Es kann somit nicht nachvollzogen werden, wie stark eine Anfrage als Angriff 
eingeordnet wird oder nicht. Dies könnte verbesert werden, indem die Wahrscheinlichkeit berechnet wird und anhand eines Schwellwertes eine ANfrage klassifziert 
werden könnte. Dies könnte bei dem Kunden service untertsützend sein, da mithilfe durch die Veränderung des SChwellwertes die \ac{fpr} beeinflusst werden könnte
und nicht so viele der Kundenanfragen abgehlehnt werden würden. \\
Die SChutzmaßnahme sollte erweitert werden, sodass die Anfragen nicht unabhängig voneinander bettrachtet werden, sondern es möglich ist die Prüfung über einen 
ganzen Chatverlauf hinweg durchzuführen. \\
Zusätzlich osllte die Prüfung der Anfrage auf externe Inhalte implementiert werden, damit eine E-Mail, welche in das Textfeld kopiert wird auch als E-Mail erkannt
wird und dementprechend als externer Inhalt geprüft wird und nicht wie es jetzt ist als nfache Nutzeranfrage. \\
Für eine bessere Übertragbarkeit sollte der Ansatz nicht nur mit weiteren \acp{llm} geprüft werden, sondern auch mit weiteren Datensätzen, auch deutschsprachige.
zusätzlich sollte eine bessere nachahmung eines Angreifers ermöglicht werden und für das Testen nicht ausschließlich statische Datensätze verwendet werden. Damit
kann die Bewertung praxisnaher werden und sich besser an die aktuelle Entwicklung anpassen. \\
Mit diesem Ausblick lässt sich abschließend festhalten, dass huetzutage die notwendigkeit von geeigneten Schuztmaßnahmen gegen ANgriffe bei Chatbots immer wichtiger
werden. Es werden immer mehr Funktionen in Chatbots integriert und immer mehr Nutzen erwartet. Jdoch ist es durch die Allgemeine Architektur eines \ac{llm} 
erschwert, das dieser SChutz ermöglicht wird bzw. die ANgriffsfläche minimiert wird. Die in dieser Arbeit vorgestellte Schutzmaßnahme zeigt vor allem die bedeutung
die vorgeschaltete Prüfungen übernehmen.